\par\Needspace{7\baselineskip}\bigskip
\begin{center}
\pdfbookmark[2]{Gauss to Olbers, Göttingen, 15 February 1817}{pallas-brief-18170215}
\textsc{Gauss} to \textsc{Olbers}, Göttingen, 15 February 1817.
\end{center}
\noindent\hspace*{\parindent}. . . . .\quad I have also begun to calculate an auxiliary table for the perturbations of Pallas, a task of about \(\frac{1}{2}\) million digits that I could never have undertaken without the active assistance of some young people, especially Mr. \textsc{Westphal}. More than half is already finished. When the table is complete, I will calculate all the oppositions once again with the utmost care and include the observations from 1802, which have so far been left out altogether. Applying probability theory already gives \textls[180]{my} determination of the mass of {\AstroSymbols ♃} the \textls[180]{same} reliability as \textsc{Bouvard}'s has according to \textsc{Laplace}'s calculation. The last opposition has not even been taken into account yet, and it will change the value only very slightly but considerably increase its reliability. The next {\AstroSymbols ☍} will do even more. Please also look for {\AstroSymbols ⚴} as soon as possible.\quad . . . . .
\par
